\chapter{2008年辽宁卷（文）}

\section{单选题}

\begin{problem}
已知集合 \(M = \{x \mid -3 < x < 1\}\)，\(N = \{x \mid x \leqslant -3\}\)，则 \(M \cup N =\)（\quad）

\choices
    {\( \varnothing \)}
    {\(\{x\mid x\geqslant -3\}\)}
    {\(\{x\mid x\geqslant 1\}\)}
    {\(\{x\mid x < 1\}\)}
\answer{D}

\begin{solution}
由集合定义，\(M=\left(-3,1\right)\)，\(N=\left(-\infty,-3\right]\). 因为 \(-3\) 属于 \(N\)，所以
\(M\cup N=\left(-\infty,1\right)\)，这正是选项 D 所表示的集合，故选 D.
\end{solution}
\end{problem}

\begin{problem}
若函数 \( y = (x + 1)(x - a) \) 为偶函数，则 \( a = \)（\quad）

\choices
    {\(-2\)}
    {\(-1\)}
    {\(1\)}
    {\(2\)}
\answer{C}

\begin{solution}
展开函数得 \(y=x^2+(1-a)x-a\). 偶函数满足 \(f(-x)=f(x)\)，因此一次项系数必须为零，即 \(1-a=0\)，解得 \(a=1\). 故选 C.
\end{solution}
\end{problem}

\begin{problem}
圆 \(x^{2} + y^{2} = 1\) 与直线 \(y = kx + 2\) 没有公共点的充要条件是（\quad）

\choices
    {\(k \in (-\sqrt{2}, \sqrt{2})\)}
    {\(k \in (-\sqrt{3}, \sqrt{3})\)}
    {\(k \in (-\infty, -\sqrt{2}) \cup (\sqrt{2}, +\infty)\)}
    {\(k \in (-\infty, -\sqrt{3}) \cup (\sqrt{3}, +\infty)\)}
\answer{B}

\begin{solution}
直线可写为 \(kx-y+2=0\)，圆心为原点，半径为 \(1\). 圆心到直线的距离为
\(\dfrac{2}{\sqrt{k^2+1}}\). 直线与圆没有公共点，当且仅当该距离大于半径，即
\(\dfrac{2}{\sqrt{k^2+1}}>1\)，所以 \(k^2<3\)，从而 \(k\in\left(-\sqrt{3},\sqrt{3}\right)\). 当 \(k^2=3\) 时直线与圆相切，不能取等号，故选 B.
\end{solution}
\end{problem}

\begin{problem}
已知 \(0 < a < 1\)，\(x = \log_a\sqrt{2} + \log_a\sqrt{3}\)，\(y = \frac{1}{2}\log_a\sqrt{5}\)，\(z = \log_a\sqrt{21} - \log_a\sqrt{3}\). 则（\quad）

\choices
    {\(x > y > z\)}
    {\(z > y > x\)}
    {\(y > x > z\)}
    {\(z > x > y\)}
\answer{C}

\begin{solution}
利用对数的运算性质，\(x=\log_a\sqrt{6}\)，\(y=\log_a\sqrt[4]{5}\)，\(z=\log_a\sqrt{7}\). 对正数比较可得
\((\sqrt{7})^2=7>6=(\sqrt{6})^2\)，且 \((\sqrt{6})^4=36>5=(\sqrt[4]{5})^4\)，所以
\(\sqrt{7}>\sqrt{6}>\sqrt[4]{5}\). 又 \(0<a<1\) 时函数 \(\log_a t\) 在正数范围内单调递减，所以 \(z<x<y\). 因而 \(y>x>z\)，故选 C.
\end{solution}
\end{problem}

\begin{problem}
已知四边形 \(ABCD\) 的三个顶点 \(A(0,2)\)，\(B(-1,-2)\)，\(C(3,1)\)，且 \(\overrightarrow{BC} = 2\overrightarrow{AD}\)，则顶点 \(D\) 的坐标为（\quad）

\choices
    {\( \left(2,\frac{7}{2}\right) \)}
    {\( \left(2,-\frac{1}{2}\right) \)}
    {\((3, 2)\)}
    {\((1,3)\)}
\answer{A}

\begin{solution}
\(\overrightarrow{BC}=(3-(-1),1-(-2))=(4,3)\). 由 \(\overrightarrow{BC}=2\overrightarrow{AD}\)，得
\(\overrightarrow{AD}=\left(2,\frac{3}{2}\right)\). 因此
\(D=A+\overrightarrow{AD}=\left(0,2\right)+\left(2,\frac{3}{2}\right)=\left(2,\frac{7}{2}\right)\)，故选 A.
\end{solution}
\end{problem}

\begin{problem}
设 \(P\) 为曲线 \(C: y = x^{2} + 2x + 3\) 上的点，且曲线 \(C\) 在点 \(P\) 处切线倾斜角的取值范围是 \(\left[0, \frac{\pi}{4}\right]\)，则点 \(P\) 横坐标的取值范围是（\quad）

\choices
    {\(\left[-1, -\frac{1}{2}\right]\)}
    {\([-1, 0]\)}
    {\([0, 1]\)}
    {\(\left[\frac{1}{2}, 1\right]\)}
\answer{A}

\begin{solution}
曲线 \(y=x^2+2x+3\) 在横坐标为 \(x\) 的点处的切线斜率为
\(y'=2x+2\). 切线倾斜角 \(\theta\in\left[0,\frac{\pi}{4}\right]\) 时，\(\tan\theta\in[0,1]\)，所以
\(0\leqslant2x+2\leqslant1\). 解得 \(-1\leqslant x\leqslant-\frac{1}{2}\)，故选 A.
\end{solution}
\end{problem}

\begin{problem}
\(4\text{ 张}\)卡片上分别写有数字 \(1\)，\(2\)，\(3\)，\(4\)，从这 \(4\text{ 张}\)卡片中随机抽取 \(2\text{ 张}\)，则取出的 \(2\text{ 张}\)卡片上的数字之和为奇数的概率为（\quad）

\choices
    {\(\frac{1}{3}\)}
    {\( \frac{1}{2} \)}
    {\( \frac{2}{3} \)}
    {\( \frac{3}{4} \)}
\answer{C}

\begin{solution}
从 \(1\)，\(2\)，\(3\)，\(4\) 中抽取 \(2\) 张的基本事件总数为 \(\mathrm C_4^2=6\). 两数之和为奇数，当且仅当抽到一张奇数卡片和一张偶数卡片，因此有 \(\mathrm C_2^1\mathrm C_2^1=4\) 个有利事件. 所求概率为
\(\dfrac{4}{6}=\dfrac{2}{3}\)，故选 C.
\end{solution}
\end{problem}

\begin{problem}
将函数 \( y = 2^{x} + 1 \) 的图象按向量 \(\bs{a}\) 平移得到函数 \( y = 2^{x + 1} \) 的图象，则（\quad）

\choices
    {\(\bs{a} = (-1, - 1)\)}
    {\(\bs{a} = (1, - 1)\)}
    {\(\bs{a} = (1,1)\)}
    {\(\bs{a} = (-1,1)\)}
\answer{A}

\begin{solution}
设平移向量为 \(\bs a=(u,v)\). 将 \(y=2^x+1\) 的图象向右、向上分别平移 \(u\)、\(v\) 后，所得图象方程为
\(y=2^{x-u}+1+v\). 要使它恒等于 \(y=2^{x+1}\)，必须有 \(-u=1\) 且 \(1+v=0\)，所以 \(u=-1\)，\(v=-1\). 因而 \(\bs a=(-1,-1)\)，故选 A.
\end{solution}
\end{problem}

\begin{problem}
已知变量 \(x\)，\(y\) 满足约束条件 \(\begin{cases}
y + x - 1 \leqslant 0,\\
y - 3x - 1 \leqslant 0,\\
y - x + 1 \geqslant 0.
\end{cases}\)，则 \(z = 2x + y\) 的最大值为（\quad）

\choices
    {\(4\)}
    {\(2\)}
    {\(1\)}
    {\(-4\)}
\answer{B}

\begin{solution}
三个约束分别化为 \(y\leqslant1-x\)、\(y\leqslant3x+1\) 和 \(y\geqslant x-1\). 可行域的三个顶点由边界直线两两相交得到：
\((-1,-2)\)，\((0,1)\)，\((1,0)\).
因为 \(z=2x+y\) 是线性函数，只需比较它在三个顶点处的值：
\(z=-4\)、\(z=1\)、\(z=2\). 最大值为 \(2\)，故选 B.
\end{solution}
\end{problem}

\begin{problem}
一生产过程有 \(4\text{ 道}\)工序，每道工序需要安排一人照看. 现从甲，乙，丙等 \(6\text{ 名}\)工人中安排 \(4\text{ 人}\)分别照看一道工序，第一道工序只能从甲，乙两工人中安排 \(1\text{ 人}\)，第四道工序只能从甲，丙两工人中安排 \(1\text{ 人}\)，则不同的安排方案共有（\quad）

\choices
    {\(24\text{ 种}\)}
    {\(36\text{ 种}\)}
    {\(48\text{ 种}\)}
    {\(72\text{ 种}\)}
\answer{B}

\begin{solution}
若第一道工序安排甲，则第四道工序只能安排丙，此时第二、三道工序从其余 \(4\) 人中有序选取，方案数为 \(\mathrm A_4^2=12\).

若第一道工序安排乙，则第四道工序可安排甲或丙. 对这两种安排，第二、三道工序均有 \(\mathrm A_4^2=12\) 种安排. 因此总方案数为
\(12+12+12=36\)，故选 B.
\end{solution}
\end{problem}

\begin{problem}
已知双曲线 \(9y^{2} - m^{2}x^{2} = 1\)（\(m > 0\)）的一个顶点到它的一条渐近线的距离为 \(\frac{1}{5}\)，则 \(m =\)（\quad）

\choices
    {\(1\)}
    {\(2\)}
    {\(3\)}
    {\(4\)}
\answer{D}

\begin{solution}
将双曲线方程化为
\(\dfrac{y^2}{(1/3)^2}-\dfrac{x^2}{(1/m)^2}=1\)，故其顶点为 \(\left(0,\frac{1}{3}\right)\) 和 \(\left(0,-\frac{1}{3}\right)\)，渐近线为 \(y=\pm\frac{m}{3}x\). 取顶点 \(\left(0,\frac{1}{3}\right)\) 和渐近线 \(mx-3y=0\)，则点到直线的距离为
\[
\frac{\left|m\cdot0-3\cdot\frac{1}{3}\right|}{\sqrt{m^2+9}}=\frac{1}{\sqrt{m^2+9}}
\]
由题设 \(\dfrac{1}{\sqrt{m^2+9}}=\dfrac{1}{5}\)，得 \(m^2+9=25\). 又 \(m>0\)，所以 \(m=4\)，故选 D.
\end{solution}
\end{problem}

\begin{problem}
在正方体 \(ABCD - A_{1}B_{1}C_{1}D_{1}\) 中，\(E\)，\(F\) 分别为棱 \(AA_{1}\)，\(CC_{1}\) 的中点，则在空间中与三条直线 \(A_{1}D_{1}\)，\(EF\)，\(CD\) 都相交的直线（\quad）

\choices
    {不存在}
    {有且只有两条}
    {有且只有三条}
    {有无数条}
\answer{D}

\begin{solution}
建立坐标系，取
\(A(0,0,0)\)，\(B(1,0,0)\)，\(C(1,1,0)\)，\(D(0,1,0)\)，\(A_1(0,0,1)\)，\(D_1(0,1,1)\). 由于 \(E\)、\(F\) 分别是 \(AA_1\)、\(CC_1\) 的中点，
\(E(0,0,\frac{1}{2})\)、\(F(1,1,\frac{1}{2})\). 因而直线 \(EF\) 的方程为
\(x=y\)，\(\ z=\frac{1}{2}\).

任取 \(t\in\mathbb R\)，取 \(X_t(0,t,1)\in A_1D_1\) 和 \(Z_t(t+1,1,0)\in CD\). 直线 \(X_tZ_t\) 与平面 \(z=\frac{1}{2}\) 的交点为
\[
Y_t=\frac{X_t+Z_t}{2}=\left(\frac{t+1}{2},\frac{t+1}{2},\frac{1}{2}\right)\in EF
\]
所以直线 \(X_tZ_t\) 同时经过 \(A_1D_1\) 上的 \(X_t\)、\(CD\) 上的 \(Z_t\) 和 \(EF\) 上的 \(Y_t\). 不同的 \(t\) 给出不同的直线，故这样的直线有无数条，选 D.
\end{solution}
\end{problem}

\section{填空题}

\begin{problem}
函数 \( y = \e^{2x + 1}\)（\(-\infty < x < +\infty\)）的反函数是\fillinblank{}.
\answer{\(y = \frac{1}{2}(\ln x - 1)\)，\(x > 0\)}

\begin{solution}
由 \(y=\e^{2x+1}\) 得 \(y>0\)，且 \(\ln y=2x+1\)，所以
\(x=\dfrac{\ln y-1}{2}\). 交换 \(x\)、\(y\) 得反函数
\(y=\dfrac{1}{2}(\ln x-1)\)，其定义域为原函数的值域 \(x>0\).
\end{solution}
\end{problem}

\begin{problem}
在体积为 \(4\sqrt{3}\pi\) 的球的表面上有 \(A\)，\(B\)，\(C\) 三点，\(AB = 1\)，\(BC = \sqrt{2}\)，\(A\)，\(C\) 两点的球面距离为 \(\frac{\sqrt{3}}{3}\pi\)，则球心到平面 \(ABC\) 的距离为\fillinblank{}.
\answer{\(\frac{3}{2}\)}

\begin{solution}
设球的半径为 \(R\). 由球的体积公式，
\(\dfrac{4}{3}\pi R^3=4\sqrt{3}\pi\)，所以 \(R^3=3\sqrt{3}\)，即 \(R=\sqrt{3}\).

设球心为 \(O\). 由 \(AC\) 的球面距离为 \(\dfrac{\sqrt{3}}{3}\pi\)，得
\(\angle AOC=\dfrac{(\sqrt{3}/3)\pi}{\sqrt{3}}=\dfrac{\pi}{3}\). 因而
\(AC=2R\sin\dfrac{\angle AOC}{2}=2\sqrt{3}\sin\dfrac{\pi}{6}=\sqrt{3}\).
又 \(AB^2+BC^2=1+2=3=AC^2\)，所以 \(\angle ABC=\dfrac{\pi}{2}\). 直角三角形 \(ABC\) 的外接圆半径为
\(\dfrac{AC}{2}=\dfrac{\sqrt{3}}{2}\).

设 \(H\) 为球心 \(O\) 到平面 \(ABC\) 的垂足. 由于 \(OA=OB=OC\)，点 \(H\) 是三角形 \(ABC\) 的外心，所以
\[
OH^2+\left(\frac{\sqrt{3}}{2}\right)^2=R^2
\]
从而 \(OH^2=3-\dfrac{3}{4}=\dfrac{9}{4}\)，故球心到平面 \(ABC\) 的距离为 \(OH=\dfrac{3}{2}\).
\end{solution}
\end{problem}

\begin{problem}
\((1 + x^{3})\left(x + \frac{1}{x^{2}}\right)^{6}\) 展开式中的常数项为\fillinblank{}.
\answer{35}

\begin{solution}
在 \(\left(x+\dfrac{1}{x^2}\right)^6\) 中，若有 \(j\) 个因子取 \(\dfrac{1}{x^2}\)，则该项的幂为
\(x^{6-3j}\)，系数为 \(\mathrm C_6^j\). 乘以 \(1+x^3\) 后，常数项有两种来源：

当取 \(1\) 时，需 \(6-3j=0\)，即 \(j=2\)，所得系数为 \(\mathrm C_6^2=15\)；

当取 \(x^3\) 时，需 \(6-3j+3=0\)，即 \(j=3\)，所得系数为 \(\mathrm C_6^3=20\).

所以常数项为 \(15+20=35\).
\end{solution}
\end{problem}

\begin{problem}
设 \(x \in \left(0, \frac{\pi}{2}\right)\)，则函数 \(y = \frac{2\sin^2 x + 1}{\sin 2x}\) 的最小值为\fillinblank{}.
\answer{\(\sqrt{3}\)}

\begin{solution}
令 \(t=\tan x\). 由于 \(x\in\left(0,\frac{\pi}{2}\right)\)，所以 \(t>0\). 利用
\(\sin^2x=\dfrac{t^2}{1+t^2}\) 和 \(\sin2x=\dfrac{2t}{1+t^2}\)，得
\[
y=\frac{2\sin^2x+1}{\sin2x}
=\frac{(3t^2+1)/(1+t^2)}{2t/(1+t^2)}
=\frac{3t+\frac{1}{t}}{2}
\]
由基本不等式，\(3t+\dfrac{1}{t}\geqslant2\sqrt{3t\cdot\dfrac{1}{t}}=2\sqrt{3}\)，等号当且仅当 \(3t=\dfrac{1}{t}\)，即 \(t=\dfrac{1}{\sqrt{3}}>0\). 该值可以取到，因此函数的最小值为 \(\sqrt{3}\).
\end{solution}
\end{problem}

\section{解答题}

\begin{problem}
在 \(\triangle ABC\) 中，内角 \(A\)，\(B\)，\(C\) 对边的边长分别是 \(a\)，\(b\)，\(c\). 已知 \(c = 2\)，\(C = \frac{\pi}{3}\).

\begin{enumerate}
\item 若 \(\triangle ABC\) 的面积等于 \(\sqrt{3}\)，求 \(a\)，\(b\)；
\item 若 \(\sin B = 2\sin A\)，求 \(\triangle ABC\) 的面积.
\end{enumerate}
\end{problem}

\begin{solution}
\begin{enumerate}
\item 由三角形面积公式，得 \(\frac{1}{2}ab\sin C=\sqrt{3}\)，所以 \(\frac{\sqrt{3}}{4}ab=\sqrt{3}\)，即 \(ab=4\). 由余弦定理，\(c^2=a^2+b^2-2ab\cos C\)，从而 \(4=a^2+b^2-ab=a^2+b^2-4\)，故 \(a^2+b^2=8\). 因此 \((a-b)^2=a^2+b^2-2ab=8-8=0\)，结合边长为正，得 \(a=b=2\).

\item 由正弦定理，\(\dfrac{\sin B}{\sin A}=\dfrac{b}{a}\). 已知 \(\sin B=2\sin A\)，所以 \(b=2a\). 再由余弦定理，\(4=a^2+b^2-2ab\cos C=a^2+4a^2-2a\cdot2a\cdot\frac{1}{2}=3a^2\)，得 \(a=\frac{2\sqrt{3}}{3}\)，\(b=\frac{4\sqrt{3}}{3}\). 因而三角形面积为 \(\frac{1}{2}ab\sin C=\frac{1}{2}\cdot\frac{8}{3}\cdot\frac{\sqrt{3}}{2}=\frac{2\sqrt{3}}{3}\).
\end{enumerate}
\end{solution}

\begin{problem}
某批发市场对某种商品的周销售量（单位：吨）进行统计，最近 \(100\text{ 周}\)的统计结果如下表所示：

\begin{center}
\begin{tabular}{c|ccc}
\hline
周销售量 & \(2\) & \(3\) & \(4\) \\
\hline
频数 & \(20\) & \(50\) & \(30\) \\
\hline
\end{tabular}
\end{center}

\begin{enumerate}
\item 根据上面统计结果，求周销售量分别为 \(2\text{ 吨}\)，\(3\text{ 吨}\)和 \(4\text{ 吨}\)的频率；
\item 若以上述频率作为概率，且各周的销售量相互独立，求
\begin{enumerate}
\item \(4\text{ 周}\)中该种商品至少有一周的销售量为 \(4\text{ 吨}\)的概率；
\item 该种商品 \(4\text{ 周}\)的销售量总和至少为 \(15\text{ 吨}\)的概率.
\end{enumerate}
\end{enumerate}
\end{problem}

\begin{solution}
\begin{enumerate}
\item 由频数除以总周数，得周销售量分别为 \(2\text{ 吨}\)，\(3\text{ 吨}\)和 \(4\text{ 吨}\)的频率为 \(\frac{20}{100}=0.2\)，\(\frac{50}{100}=0.5\) 和 \(\frac{30}{100}=0.3\).

\item
\begin{enumerate}
\item 每周销售量为 \(4\text{ 吨}\)的概率为 \(0.3\)，每周不是 \(4\text{ 吨}\)的概率为 \(1-0.3=0.7\). 由各周销售量相互独立，\(4\) 周中至少有一周销售量为 \(4\text{ 吨}\)的概率为 \(1-0.7^4=1-0.2401=0.7599\).

\item 四周销售量的总和至少为 \(15\text{ 吨}\)，而每周销售量至多为 \(4\text{ 吨}\)，所以只有两种情形：四周均为 \(4\text{ 吨}\)，或恰有一周为 \(3\text{ 吨}\)且其余三周为 \(4\text{ 吨}\). 因而所求概率为 \(0.3^4+\mathrm C_4^1\cdot0.5\cdot0.3^3=0.0081+0.054=0.0621\).
\end{enumerate}
\end{enumerate}
\end{solution}

\begin{problem}
如图，在棱长为 \(1\) 的正方体 \(ABCD - A'B'C'D'\) 中，\(AP = BQ = b\)（\(0 < b < 1\)），截面 \(PQEF \parallel A'D\)，截面 \(PQGH \parallel AD'\).

\begin{enumerate}
\item 证明：平面 \(PQEF\) 和平面 \(PQGH\) 互相垂直；
\item 证明：截面 \(PQEF\) 和截面 \(PQGH\) 面积之和是定值，并求出这个值；
\item 若 \( b = \frac{1}{2} \)，求 \(D'E\) 与平面 \(PQEF\) 所成角的正弦值.
\end{enumerate}

\bitmapfigure[max width=4.6cm]{img/2008/liaoning_liberal/q19_fig1.png}
\end{problem}

\begin{solution}
\begin{enumerate}
\item 建立如图所示的空间直角坐标系，取
\(A(0,0,0)\)，\(B(1,0,0)\)，\(C(1,1,0)\)，\(D(0,1,0)\)，以及
\(A'(0,0,1)\)，\(B'(1,0,1)\)，\(C'(1,1,1)\)，\(D'(0,1,1)\). 由 \(AP=BQ=b\)，得 \(P(0,0,b)\)，\(Q(1,0,b)\)，所以 \(\overrightarrow{PQ}=(1,0,0)\). 又 \(\overrightarrow{A'D}=(0,1,-1)\)，\(\overrightarrow{AD'}=(0,1,1)\).

设平面 \(PQEF\) 和平面 \(PQGH\) 分别为 \(\Pi_1\) 和 \(\Pi_2\). 由平行关系，\(\Pi_1\) 的一个法向量可取 \(\bs n_1=\overrightarrow{PQ}\times\overrightarrow{A'D}=(0,1,1)\)，\(\Pi_2\) 的一个法向量可取 \(\bs n_2=\overrightarrow{PQ}\times\overrightarrow{AD'}=(0,-1,1)\). 因为 \(\bs n_1\cdot\bs n_2=0\)，所以 \(\Pi_1\perp\Pi_2\)，即平面 \(PQEF\) 和平面 \(PQGH\) 互相垂直.

\item 平面 \(\Pi_1\) 过点 \(P\)，且包含方向向量 \(\overrightarrow{PQ}\) 和 \(\overrightarrow{A'D}\)，故其方程为 \(y+z=b\). 它与底面 \(z=0\) 的交线为 \(EF\)，其中 \(E(1,b,0)\)，\(F(0,b,0)\). 因而 \(PQ=1\)，\(QE=\sqrt{b^2+b^2}=b\sqrt{2}\)，所以 \(S_{PQEF}=b\sqrt{2}\).

平面 \(\Pi_2\) 的方程为 \(z-y=b\). 它与顶面 \(z=1\) 的交线为 \(GH\)，其中 \(G(1,1-b,1)\)，\(H(0,1-b,1)\). 因而 \(PQ=1\)，\(QG=\sqrt{(1-b)^2+(1-b)^2}=(1-b)\sqrt{2}\)，所以 \(S_{PQGH}=(1-b)\sqrt{2}\). 因此两截面面积之和为 \(S_{PQEF}+S_{PQGH}=b\sqrt{2}+(1-b)\sqrt{2}=\sqrt{2}\)，是定值.

\item 当 \(b=\frac{1}{2}\) 时，\(E(1,\frac{1}{2},0)\)，所以 \(\overrightarrow{D'E}=(1,-\frac{1}{2},-1)\). 平面 \(PQEF\) 的法向量为 \(\bs n_1=(0,1,1)\). 设直线 \(D'E\) 与平面 \(PQEF\) 所成的角为 \(\theta\)，则
\[
\sin\theta=\frac{\lvert\overrightarrow{D'E}\cdot\bs n_1\rvert}{\lvert\overrightarrow{D'E}\rvert\,\lvert\bs n_1\rvert}=\frac{\left\lvert-\frac{1}{2}-1\right\rvert}{\sqrt{1+\frac{1}{4}+1}\sqrt{2}}=\frac{\frac{3}{2}}{\frac{3}{2}\sqrt{2}}=\frac{\sqrt{2}}{2}
\]
\end{enumerate}
\end{solution}

\begin{problem}
在数列 \(\{a_{n}\}\)，\(\{b_{n}\}\) 是各项均为正数的等比数列，设 \(c_{n} = \frac{b_{n}}{a_{n}}\)（\(n \in \N^{*}\)）.

\begin{enumerate}
\item 数列 \(\{c_{n}\}\) 是否为等比数列？证明你的结论；
\item 设数列 \(\{\ln a_n\}\)，\(\{\ln b_n\}\) 的前 \(n\) 项和分别为 \(S_n\)，\(T_n\). 若 \(a_1 = 2\)，\(\frac{S_n}{T_n} = \frac{n}{2n + 1}\)，求数列 \(\{c_n\}\) 的前 \(n\) 项和.
\end{enumerate}
\end{problem}

\begin{solution}
\begin{enumerate}
\item 设 \(\{a_n\}\) 和 \(\{b_n\}\) 的公比分别为 \(r\) 和 \(s\). 由于两数列各项均为正数，\(r>0\)，\(s>0\). 于是
\(\dfrac{c_{n+1}}{c_n}=\dfrac{b_{n+1}/a_{n+1}}{b_n/a_n}=\dfrac{b_{n+1}}{b_n}\cdot\dfrac{a_n}{a_{n+1}}=\dfrac{s}{r}\)，它与 \(n\) 无关，所以 \(\{c_n\}\) 是等比数列，公比为 \(\dfrac{s}{r}\).

\item 由 \(a_1=2\)，可设 \(a_n=2r^{n-1}\)，\(b_n=b_1s^{n-1}\). 记 \(L=\ln2\)，\(\alpha=\ln r\)，\(\beta=\ln s\)，\(u=\ln b_1\). 则
\(S_n=nL+\frac{n(n-1)}{2}\alpha\)，\(T_n=nu+\frac{n(n-1)}{2}\beta\).
由题设关系分别取 \(n=1\)，\(2\)，\(3\)，得到
\(\frac{L}{u}=\frac{1}{3}\)，\(\frac{2L+\alpha}{2u+\beta}=\frac{2}{5}\)，\(\frac{3L+3\alpha}{3u+3\beta}=\frac{3}{7}\).
因为 \(u=3L\)，后两式分别化为
\(5\alpha-2\beta=2L\)，\(7\alpha-3\beta=2L\).
两式相减得 \(2\alpha-\beta=0\)，即 \(\beta=2\alpha\). 代入第一式得 \(\alpha=2L\)，从而 \(\beta=4L\). 因此 \(r=4\)，\(s=16\)，且 \(b_1=8\).

于是 \(c_n=\dfrac{8\cdot16^{n-1}}{2\cdot4^{n-1}}=4^n\)，故数列 \(\{c_n\}\) 的前 \(n\) 项和为
\[
\sum_{k=1}^{n}c_k=4+4^2+\cdots+4^n=\frac{4(4^n-1)}{4-1}=\frac{4^{n+1}-4}{3}
\]
\end{enumerate}
\end{solution}

\begin{problem}
在直角坐标系 \(xOy\) 中，点 \(P\) 到两点 \((0, -\sqrt{3})\)，\((0, \sqrt{3})\) 的距离之和为 \(4\)，设点 \(P\) 的轨迹为 \(C\).

\begin{enumerate}
\item 写出 \(C\) 的方程；
\item 设直线 \( y = kx + 1 \) 与 \(C\) 交于 \(A\)，\(B\) 两点，\( k \) 为何值时 \( \overrightarrow{OA} \perp \overrightarrow{OB}\)？此时 \( \lvert AB\rvert \) 的值是多少？
\end{enumerate}
\end{problem}

\begin{solution}
\begin{enumerate}
\item 椭圆的两个焦点为 \((0,-\sqrt{3})\) 和 \((0,\sqrt{3})\)，所以焦距的一半为 \(c=\sqrt{3}\). 由点 \(P\) 到两焦点的距离之和为 \(4\)，得长轴长 \(2a=4\)，即 \(a=2\). 因而 \(b^2=a^2-c^2=4-3=1\)，且长轴在 \(y\) 轴上，所以 \(C\) 的方程为
\[
\frac{x^2}{b^2}+\frac{y^2}{a^2}=1,
\quad\text{即}\quad
\frac{x^2}{1}+\frac{y^2}{4}=1
\]

\item 设直线与椭圆的交点横坐标分别为 \(x_1\)，\(x_2\). 将 \(y=kx+1\) 代入椭圆方程，得
\(\frac{x^2}{1}+\frac{(kx+1)^2}{4}=1\)，\((k^2+4)x^2+2kx-3=0\).

由韦达定理，
\(x_1+x_2=-\frac{2k}{k^2+4}\)，\(x_1x_2=-\frac{3}{k^2+4}\).

由于 \(A(x_1,kx_1+1)\)，\(B(x_2,kx_2+1)\)，且 \(\overrightarrow{OA}\perp\overrightarrow{OB}\)，所以
\[
0=\overrightarrow{OA}\cdot\overrightarrow{OB}=x_1x_2+(kx_1+1)(kx_2+1)
 = (1+k^2)x_1x_2+k(x_1+x_2)+1
 =\frac{1-4k^2}{k^2+4}
\]
因此 \(k^2=\frac{1}{4}\)，即 \(k=\frac{1}{2}\) 或 \(k=-\frac{1}{2}\).

又因为两交点在同一直线上，
\[
\lvert AB\rvert=\sqrt{1+k^2}\,\lvert x_1-x_2\rvert
\]
当 \(k^2=\frac{1}{4}\) 时，二次方程的判别式为
\[
\Delta=(2k)^2-4(k^2+4)(-3)=16k^2+48=52
\]
所以
\[
\lvert x_1-x_2\rvert=\frac{\sqrt{\Delta}}{k^2+4}
=\frac{2\sqrt{13}}{17/4}
=\frac{8\sqrt{13}}{17}
\]
从而
\[
\lvert AB\rvert=\sqrt{\frac{5}{4}}\cdot\frac{8\sqrt{13}}{17}
=\frac{4\sqrt{65}}{17}
\]
\end{enumerate}
\end{solution}

\begin{problem}
设函数 \( f(x) = ax^3 + bx^2 - 3a^2x + 1 \)（\(a\)，\(b \in \R\)）在 \(x = x_1\)，\(x = x_2\) 处取得极值，且 \(\lvert x_1 - x_2\rvert = 2\).

\begin{enumerate}
\item 若 \(a = 1\)，求 \(b\) 的值，并求 \(f(x)\) 的单调区间；
\item 若 \(a > 0\)，求 \(b\) 的取值范围.
\end{enumerate}
\end{problem}

\begin{solution}
\begin{enumerate}
\item 对函数求导，得 \(f'(x)=3ax^2+2bx-3a^2\). 因为在两个不同点处取得极值，所以 \(a\neq0\)，且 \(x_1\)，\(x_2\) 是方程 \(3ax^2+2bx-3a^2=0\) 的两个不等实根. 由韦达定理，
\(x_1+x_2=-\frac{2b}{3a}\)，\(x_1x_2=-a\).
于是
\[
(x_1-x_2)^2=(x_1+x_2)^2-4x_1x_2=\frac{4b^2}{9a^2}+4a
\]
结合 \(\lvert x_1-x_2\rvert=2\)，得 \(\frac{4b^2}{9a^2}+4a=4\)，即 \(b^2=9a^2(1-a)\). 当 \(a=1\) 时，\(b=0\)，此时 \(f'(x)=3x^2-3=3(x-1)(x+1)\). 因而 \(f(x)\) 在 \((-\infty,-1)\) 和 \((1,+\infty)\) 上单调递增，在 \((-1,1)\) 上单调递减.

\item 当 \(a>0\) 时，由 \(b^2=9a^2(1-a)\geqslant0\)，得 \(0<a\leqslant1\). 在此范围内，
\[
b^2=9a^2(1-a)
\]
令 \(h(a)=a^2(1-a)\)，则 \(h'(a)=a(2-3a)\)，所以 \(h(a)\) 在 \(a=\frac{2}{3}\) 处取得最大值 \(\frac{4}{27}\). 因此 \(b^2\leqslant9\cdot\frac{4}{27}=\frac{4}{3}\)，即 \(\lvert b\rvert\leqslant\frac{2\sqrt{3}}{3}\). 反过来，函数 \(3a\sqrt{1-a}\) 在 \(\left[\frac{2}{3},1\right]\) 上连续，取值从 \(\frac{2\sqrt{3}}{3}\) 变化到 \(0\)，故上述范围内的每个 \(\lvert b\rvert\) 均可取得，最终
\[
-\frac{2\sqrt{3}}{3}\leqslant b\leqslant\frac{2\sqrt{3}}{3}
\]
\end{enumerate}
\end{solution}
